\documentclass[11pt]{article}
\usepackage[a4paper,margin=25mm]{geometry}
\usepackage{amsmath,amssymb,amsthm,booktabs,hyperref}
\emergencystretch=1em
\newtheorem{theorem}{Theorem}
\title{Complete finite exclusions for integer equal-area rectangle partitions}
\author{Private recovery manuscript for L018}\date{12 October 2026}
\begin{document}\maketitle
This paper recovers the principal finite conclusions of Garc\'ia-Col\'in,
Leemans, M\"u\ss ig and Rold\'an~\cite{source}, using the complete companion
\emph{Boundary constraints for equal-area rectangle dissections} and the whole
finite programs and receipts specified below. It also records corrections to
particular examples and implementation conventions. It does not assert that
every literal source statement is true. The source's question about an integer
perfect Mondrian partition of any size remains unanswered here.

\begin{theorem}[Complete bounded conclusions]
There is no axis-parallel partition of an integer-sided square of side at most
$1000$, or of an integer-sided rectangle whose two sides are at most $400$,
into at least two pairwise noncongruent integer-sided rectangles of equal area.
Congruence includes a quarter-turn.
\end{theorem}
\begin{proof}
By the companion's seven-piece theorem, any such partition has $r\ge7$ pieces.
Their common integer area is $A=WH/r$, so $r$ is a divisor of $WH$. Each shape
is an unordered pair $a\le b$ with $ab=A$ and a fitting orientation. Interchanging
the two axes if needed, take $W\ge H$; the exact fitting condition is $a\le H$
and $b\le W$. There must be at least $r$ such shapes. The exact divisor and
factor-pair enumeration in \texttt{enumerate\_candidates.py} gives
\[
\begin{array}{c|r|r}
\text{range}&\text{containers with an arithmetic candidate}&(W,H,r)\text{ cases}\\\hline
W=H\le1000&49&103\\
H\le W\le400&6431&11126
\end{array}
\]
All other parameter choices are excluded by these necessary conditions.
The complete roster is stored explicitly, and its checker reconstructs it
from every container and every divisor, rather than trusting the displayed counts.

Every listed case has a completed exhaustive-search exclusion. The companion
proves that the finite skyline search follows any actual tiling and proves all
prunings and symmetry reductions used here. The exact per-case route in
\texttt{search\_coverage.json} names the whole program, successful receipt,
record index, and actual node count. It contains exactly the roster and no
unresolved case. Earlier bounded runs marked \texttt{UNKNOWN} are retained as
history, but no such record supplies an exclusion. This completed computation
is an explicit finite proof premise, with ordinary compiler and machine-execution
trust; it is not a kernel-checked proof or merely a sample of placements.
\end{proof}

\paragraph{The two largest square cases.}
The general exact search closes every square case except $(840,840,20)$ and
$(840,840,21)$. For these, the two complete specialized programs add exactly the
companion's proved parity reductions. In the $21$-piece case, all $21$ shapes
must be used. Four have an odd side; the normalized even partners are $7,5,3,1$.
The alternating sums force their odd dimensions onto a common axis, their
interval endpoints give six relative-height patterns, and their starting
horizontal parities separate $7+1$ from $5+3$. In the $20$-piece case, exactly
eight of the nine odd-sided shapes must be selected. Literal signed-subset
arithmetic gives $44$ configurations, or $26$ after choosing the axis with at
least four odd dimensions.

Both specialized programs retain the least-index bottom-left corner convention,
using the axis-preserving reflections. They \emph{drop} the independent
bottom-left width-at-most-height convention, since the quarter-turn has already
been used to normalize the odd dimensions. All corner comparisons are non-strict.
The two completed runs respectively explore $77\,756\,396$ and $59\,729\,842$
states and return \texttt{INFEASIBLE}. Their observed times, about $41.8$ and
$31.9$ seconds on this machine, are historical run measurements and are not
complexity bounds. The distinct signed-sum/endpoint checker reconstructs the
six and $44/26$ lists and their reflection/rotation closure independently.

\paragraph{Whole finite dependency closure.}
The required executable sources are the complete divisor-roster generator,
five complete C++ search versions, and the complete receipt-coverage assembler.
The five versions preserve the exact bytes that produced the selected receipts:
original lowest-column search, added side-subset pruning, most-constrained valley
order, and the two specialized parity searches. The whole roster and all ten
untrimmed receipt files are required mathematical data. The current assembly,
metadata, binding and render driver is operational. The independent small-cell
oracle, congruence-relaxed positive tests, source-example checker, and coloring
and specialization checks are bounded validation; they do not replace the
completed finite searches. Earlier exploratory CP-SAT models are not premises.
No new Lean theorem is claimed, and no external source census is assumed.

\paragraph{Source examples and corrected perimeter convention.}
The first square with enough shapes is $84$ with $r=7$ and common area $1008$.
Its unordered shape list is
\[
 (12,84),(14,72),(16,63),(18,56),(21,48),(24,42),(28,36).
\]
Several products in the printed list on physical page 4 differ from $1008$;
the list above is the exact factor-pair replacement. The bottom-side example
$12+72=84$ remains an example of partial boundary coverage.

The statement that a possible perimeter first occurs at $360$ needs an extra
no-spanning convention. An actual five-piece equal-area, noncongruent perimeter
in the $84$-square is given by the following $(x,y,w,h)$ coordinates:
\[
\begin{gathered}
(72,0,12,84),\ (0,0,16,63),\ (16,0,56,18),\\
(0,63,48,21),\ (48,42,24,42).
\end{gathered}
\]
These interiors are disjoint and cover all four sides; the first piece spans
the full height and occupies two corners. This is not a perfect partition.
Removing its spanning strip would leave a six-piece perfect partition of a
rectangle if a seven-piece completion existed, contrary to the companion's
seven-piece theorem. A correct general perimeter test allows opposite sides
to share a spanning piece and repeated corner identities, or explicitly pays
for a valid recursive exclusion. One cannot silently require four distinct
corner pieces in a search for arbitrary partitions.

For the source's $360$-square example, the seven perimeter pieces have coordinates
\[
\begin{gathered}
(0,0,240,45),\ (240,0,120,90),\ (0,180,60,180),\\
(60,252,100,108),\ (160,306,200,54),\\
(0,45,80,135),\ (310,90,50,216).
\end{gathered}
\]
Each area is $10800$. The five unused shapes are
$(36,300),(40,270),(48,225),(72,150),(75,144)$.
The printed assertion that their dimensions cannot sum to $216$ is false:
$72+144=216$ uses two distinct listed shapes. This arithmetic possibility is not
a tiling. The other printed arithmetic equality $40+48+72=160$ is correct,
but a $270$-high piece above the bottom gap at height $45$ reaches $315$ and
intersects the top perimeter, whose lower edges are all at most $306$.

The hole bound correctly excludes this perimeter regardless of that mistaken
subset sum. Complete left and right boundary coverings occupy vertical strips
of widths equal to their respective minimum widths. The corresponding top and
bottom strips are also occupied. Every piece not on the perimeter must therefore
fit in the bounding box
\[
 (360-60-50)\mathbin{\times}(360-45-54)=250\mathbin{\times}261.
\]
Neither orientation of the required $36\mathbin{\times}300$ piece fits.
This is a containing box, not an assertion that the interior itself is rectangular.
All displayed coordinates, areas, disjointness, side coverings and arithmetic
witnesses are checked in \texttt{source\_examples.py}.

\paragraph{Scope of method reports and earlier results.}
The stronger complete exclusions here imply the source's numerical lower ranges
$360$ and $420$ without equating the independent search with the historical
sequence of side, perimeter, gap and hole implementations. Its eleven-side
filtered list and runtime heatmap are retained as reports about that implementation,
not borrowed premises of this proof or newly reproduced method counts. The
source's current public Julia revision is pinned separately. The example
$W=60,H=14,r=7,A=120$ has seven fitting unordered shapes but eight fitting
orientations, so half the filtered orientation count is not the shape count and
reversal of that filtered list is not generally the quarter-turn map. This audits
that exact current revision; no claim about the historical execution revision
or a counterexample to the complete bounded conclusions follows.

The source's prior real dissections, earlier minimum-defect computations through
$65$, constructions through $96$, density estimates and stochastic-process links
remain expressly attributed background, not unproved premises of either bounded
theorem here. The elementary defect bounds follow by two full-height strips:
widths $(n-1)/2,(n+1)/2$ give defect $n$ for odd $n\ge3$; widths $n/2-1,n/2+1$
give defect $2n$ for even $n\ge4$. The nontrivial-partition convention excludes
$n=1,2$ from these constructions. The general integer-existence question and the
stated logarithmic defect-bound conjecture remain unanswered, with no new
question, formalization or compression credit.

\begin{thebibliography}{1}
\bibitem{source} N. Garc\'ia-Col\'in, D. Leemans, M. M\"u\ss ig and \'E. Rold\'an,
\emph{There is no perfect Mondrian partition for squares of side lengths less
than 1001}, arXiv:2311.02385v1, eleven-page fixed source edition (2023).
\end{thebibliography}\end{document}
